\subproblem{Айнымалы (?) ток}{3.5}

\pic7R{0pt}{0.3}{figures/1b.1}

Оң жақтағы суретте көрсетілген тізбекте сыйымдылығы $C=\qty{4.7}{\mF}$ конденсатор мен индуктивтілігі $L=\qty{5.4}{\H}$ катушка тізбектей жалғанып, $K$ кілтіне қосылған. Кілт екі нұсқада да $U_0=\qty{10}{\V}$ тұрақты кернеу көзіне өте қысқа $\tau=\qty5{\ms}$ уақытқа қосылады, содан кейін $T=\qty{250}{\ms}$ уақыт бойы ($T\gg\tau$ деп есептеңіз) қысқа тұйықталуға, яғни ең шеткі сол жақ тік күйге ауыстырылады, содан кейін қайтадан кернеу көзіне қосылады және т.с.с. $t=0$ сәтінде екі жағдайда да тізбекте ток болған жоқ.

\newcounter{taskstar}
\setcounter{taskstar}{0}
\begin{list}{\textbf{(\alph{taskstar})}\ }{
\usecounter{taskstar}
\setlength{\leftmargin}{18pt}%
\setlength{\labelwidth}{0pt}%
\setlength{\labelsep}{0pt}%
\setlength{\itemsep}{6pt}%
\setlength{\topsep}{6pt}}
\item Конденсатор зарядталмаған, ал кілт шеткі сол жақ тік күйде болды делік. $t=0$ сәтінде кілт $U_0$-ге қосылып, жоғарыда сипатталғандай ауысып тұрды. $t=3T+\tau$ уақыт мезетіндегі ($t=3T$, $K$ кілті шеткі сол жақ күйге ауысқаннан кейінгі мезет) жүйенің энергиясы қандай болады? $0\leq t \leq 3T+\tau$ уақыт аралығындағы конденсатордың максималды заряды $q_0$ неге тең болды?
\item Айталық, $t=0$ мезетінде конденсаторда белгісіз $q_*>0$ заряды болды (полярлығы суретте көрсетілген), ал кілт шеткі сол жақ тік күйде қозғалмай тұрды. Содан кейін, қандай да бір $t=t_*>0$ мезетінде кілт $U_0$ көзіне қосылып, жоғарыда сипатталғандай ауыстырыла бастады. $K$ кілтін ауыстырып қосқан кезде жүйенің толық энергиясы өзгеріссіз қалғаны белгілі болды. Бұл $q_*$ және $t_*$-ның қандай мәндерінде мүмкін?
\end{list}
